% Ad Atticum 10.10 (ad-atticum-10-10) --- English translation
% Translated by Claude (Anthropic) from the Latin (Perseus).
% Style guide: STYLE.md.

\ciceroLetterOpener{CICERO ATTICO salutem}

\ciceroSection{1} How blind I have been not to have seen this before! I have sent you Antony's letter. I had written to him repeatedly that I was contemplating nothing against Caesar's interests, that I bore in mind my son-in-law, that I bore in mind our friendship, that I could, if I felt otherwise, have been with Pompey, but that for myself, because I was being made to run about with lictors against my will, I wished to be elsewhere --- and that even this much I had not yet quite settled. See, then, how he replies to this --- and how exhortingly \textit{[Greek: parainetik\=os]}:

\ciceroSection{2} ``Your view, how true it is. For the man who wants to keep to the middle stays put in his own country; the man who sets out is taken to be passing judgement on one side or the other. But I am not the man who ought to be deciding by what right one person sets out and another does not; Caesar has laid this part on me, that I should not let anyone at all leave Italy. So it is of little use that I approve your reasoning, if all the same I can grant you nothing on my own. I think you should write to Caesar and ask this of him. I have no doubt you will obtain it, especially since you undertake that you will be mindful of our friendship.''

\ciceroSection{3} There you have a Spartan scytale \textit{[Greek: skytal\=en Lak\=onik\=en]}. I shall certainly receive the fellow. He was due to arrive on the fifth before the Nones in the evening, that is, today. So tomorrow perhaps he will come on to me. I shall make the attempt --- \dag{}audeam\dag{} --- to show no haste; to say I will send to Caesar. I shall act in secret; I shall conceal myself with the fewest possible companions somewhere; \dag{}carti\dag{} I shall fly off from here, much against their will, and would that it might be to Curio! Mark what I am telling you \textit{[Greek: sunes ho toi leg\=o]}. A great pain has come over me. Something worthy of us will be brought off.

\ciceroSection{4} Your strangury \textit{[Greek: dysouria]} grieves me very much. Treat it, please, while it is still at the beginning \textit{[Greek: arch\=e]}. About the Massiliotes your letter was welcome to me. I beg you to let me know whatever you hear. I would like to have Ocella with me, if I could do it openly, the arrangement I had managed through Curio. Here I am waiting for Servius; for his wife and his son keep at me, and I think there is need of it.

\ciceroSection{5} And meanwhile this fellow is carrying Cytheris around with him in an open litter, like a second wife. And besides her, seven linked litters of mistresses --- or boyfriends. See what a shameful death we are perishing by, and doubt, if you can, that he, whether he comes back beaten or victorious, will be making slaughter. But I --- even in a little skiff, if no ship is to be had --- will tear myself out of these traitors' way. But I shall write more when I have met with him.

\ciceroSection{6} Our young man --- I cannot help loving him, but I understand plainly enough that I am not loved by him. I have never seen anything so without a formed character \textit{[Greek: an\=ethopoi\=eton]}, so turned away from his own people, so cogitating who knows what. Oh, the incredible weight of troubles! But it shall be my care, and is, that he be steered. For there is a remarkable wit there; his character must be tended to \textit{[Greek: \=ethous epimel\=eteon]}.
